\section{Conclusion}
\label{sec:conclusion}

We presented a systematic empirical comparison of classical error correcting
codes and LLM-based correction on the Binary Symmetric Channel, using metrics
spanning bit-level (BER) to semantic-level (WER, BLEU) quality.

Our key findings are:
\begin{enumerate}
  \item \textbf{ECC dominates at low noise} ($p < 0.10$): algebraic correction
        is near-exact within the error-correction radius and incurs no
        hallucination risk.
  \item \textbf{LLMs provide semantic recovery at high noise} ($p \geq 0.15$):
        linguistic priors allow recovery of intended meaning even when
        bit-level correction has failed.
  \item \textbf{The combined ECC+LLM pipeline is strictly better} across all
        tested $p$ values, confirming complementarity.
  \item \textbf{Shannon's bound is respected}: no method achieves information
        rates above $C(p)$, validating the simulation framework.
\end{enumerate}

These results have practical implications for systems where text is transmitted
over noisy channels: LLM post-processing on top of standard ECC is a low-cost
enhancement that can meaningfully improve perceived quality at the semantic
level, particularly for high-noise regimes where ECC alone is insufficient.

\textbf{Future work} includes: (i) evaluating learned neural channel decoders
\cite{farsad2018deep} in the same framework; (ii) investigating fine-tuned LLMs
on synthetic BSC-corrupted data; (iii) extending to other channel models
(burst errors, erasure channels); (iv) characterizing the information-theoretic
role of the LLM source prior formally.
